\section*{Conclusion}
\addcontentsline{toc}{section}{CONCLUSION}

This thesis investigated the problem of AI-assisted directional
prediction of thermoelastic wave propagation in geological media. The
work combined a physical and geological foundation, a coupled
thermoelastic mathematical formulation in two-dimensional plane
strain, a controlled COMSOL Multiphysics dataset generation pipeline,
four neural-surrogate routes (PINN, FNO, MeshGraphNet and a
Transformer-style operator) wired through a single normalised
prediction contract, and an end-to-end software prototype that
exposes the comparison through a backend orchestrator and a web
frontend.

The theoretical part (Chapters~2 and~3) established the physical
and mathematical setting. Rocks were treated as locally homogeneous
isotropic linear thermoelastic media with effective parameters drawn
from the canonical material table of the thesis. The coupled
heat--momentum system, the thermal expansion coefficient $\alpha$, and
the thermoelastic coupling $\gamma = 3K\alpha$ provided the
constitutive backbone reused by every later chapter. Directional
quantities were defined through the source--probe geometry without
requiring a full three-dimensional field solution.

The methodological and implementation part (Chapters~4--10) described
the comparative prediction methodology, the COMSOL setup that
produced the per-rock 2D datasets, the system implementation, and the
four surrogate architectures. A v2 API contract was frozen and implemented
end-to-end: every request now collapses to four user-facing fields
(model, medium identifier, source/probe geometry, observation time),
and every model service returns a normalised response with the same
physical layout (temperature, displacement components, displacement
magnitude, directional geometry, predicted travel time, plus
diagnostics). Training-data invariants (reference temperature
$273.15$~K, source temperature $1500$~K, domain $1\,\text{m}\times
1\,\text{m}$, no source frequency) are locked server-side so clients
cannot push the trained models out of distribution.

The empirical part (Chapter~12) compared the four routes on a balanced
sweep of $40$ scenarios per material across four rocks (sandstone,
limestone, basalt, granite) under one experimental protocol. The
measurement script used for the accuracy and latency comparison is
provided in Appendix~\ref{app:accuracy-latency}. The main quantitative
findings are:

\begin{itemize}
\item
PINN is the only route whose predicted travel time agrees with the
analytical body-wave reference $t_{\mathrm{gt}} = \lVert x_p - x_s
\rVert / V_p$ in every scenario (median relative error $2\,\%$,
$100\,\%$ of scenarios within the $10\,\%$ band).
\item
The Transformer-style operator is the fastest route by per-call
latency ($\sim 1.3$~ms median) and is the most material-sensitive on
the temperature output, but its absolute travel-time prediction is
already outside the $10\,\%$ band (median relative error $\sim 29\,\%$).
\item
MeshGraphNet runs the trained rollout in the current deployment, but
its travel-time scale is dominated by the rollout horizon rather than
by the requested observation time. It is recommended as a graph-based
surrogate candidate, not as a standalone physical predictor.
\item
The Fourier neural operator passes the structural sanity tests (it
returns a value of the right type for every request) but its
displacement and temperature magnitudes in the current iteration are
many orders of magnitude away from the physically meaningful range.
This is traced to output denormalisation against
$\sigma_T = 0.017$~K training statistics; the implementation task is
to retrain with stable normalisation rather than to discard the
architecture.
\item
Only PINN produces a smooth time-dependent trace at the probe.
FNO, MeshGraphNet and the Transformer operator behave as fixed-horizon
operators on the trained scenarios: their predictions are
approximately constant across the requested observation time. This is
an architectural property surfaced honestly by the comparison
framework, not a plotting artefact.
\end{itemize}

Taken together, these findings answer the central question of the
thesis: \emph{can a unified AI-assisted prediction contract
meaningfully compare physics-informed and data-driven surrogates for
directional thermoelastic prediction in geological media?} The
answer is positive. The v2 contract holds the four routes to the same
input description and the same output schema, so the differences seen
in the results are properties of the architectures, not of the
interface. Within this comparison, the physics-informed route is the
most accurate for the body-wave reference, the operator-learning
routes are faster and more flexible, and the gaps between them are
well defined.

The work has several limitations that should be made explicit at the
defence. The training corpus is a controlled COMSOL dataset, not a
field-validated geological dataset, so the absolute numerical values
should be read as comparative indicators rather than as engineering
predictions. The COMSOL setup itself is a simplified rod-in-rock
geometry, not a full subsurface model. The catalog uses midpoint
literature parameters, not site-measured values. Six of the ten
materials in the v2 catalog have no thermal expansion coefficient in
the source CSV and are therefore gated off from thermoelastic
prediction by design. The FNO calibration problem identified in
Chapter~12 requires a retraining pass that is out of scope for the
present iteration. MeshGraphNet was loaded against a 2D dataset
padded to match a 3D checkpoint; absolute magnitudes there inherit
the cost of that mismatch.

Directions for future work follow naturally from these limitations.
Retrain FNO with checkpoint-consistent normalisation. Train
MeshGraphNet directly on the 2D rod datasets used in this thesis
rather than reusing the 3D checkpoint. Extend the catalog with
laboratory-measured thermal expansion coefficients so that more rocks
become thermoelastic-supported. Add a low-cost analytical Green's
tensor baseline alongside the analytical $r/V_p$ reference for a
richer accuracy comparison. Once the four routes converge on the
same calibrated output range, the comparison can move from
\emph{which route is physically consistent} to \emph{which route
generalises across geological materials it has not seen in training}.

The practical software prototype implemented in this thesis is not a
field-validated thermoelastic simulator and was not designed to
replace classical numerical solvers. It is a research MVP that
demonstrates a normalised contract, a working four-route comparison,
and an honest visualisation of where physics-informed and data-driven
surrogates currently differ for directional thermoelastic prediction
in geological media. Given the limitations described above and the future work outlined, the prototype provides a reproducible basis for the ongoing research of AI-assisted prediction in geophysics
